
\begin{corollary}[Every peer emits a fresh self descriptor quickly]
\label{cor:c2successfulsourcecover}
Let $\zeta_r$ be the first event at which peer $r$
initiates a nonempty-cache exchange. Uniformly over
all weakly connected states,
\[
 \mathbb E_x\zeta_r\le8n,\qquad
 \Pr_x\{\zeta_r>16nk\}\le2^{-k}.
\]
Consequently
\[
 \Pr_x\!\left\{\max_r\zeta_r>
 16n\left\lceil\log_2(n/\varepsilon)\right\rceil
 \right\}\le\varepsilon.
\]
This means every peer has sent its fresh self
descriptor; it does not mean all old descriptor
copies have disappeared or samples are uniform.
\end{corollary}
\begin{proof}
If $C_r$ is nonempty, passive exchanges cannot
empty it, because at least one valid fresh self
identifier is received and protected. Until $r$
initiates, it therefore stays nonempty.
The next independent initiator clock of $r$
has mean waiting time $n$ events and produces
the desired exchange.

If $C_r$ is empty, weak connectivity supplies
some other cache containing $r$.
Apply Theorem~\ref{thm:c2readycontact} to the cut
$\{r\},V\setminus\{r\}$. Its first exchange has
mean waiting time at most $7n$ and must arrive
at $r$, since an empty-cache initiator is a
no-op. It makes $C_r$ nonempty.
The preceding argument then costs at most
another $n$ events, giving $\mathbb E\zeta_r\le8n$.
Markov's inequality and restart before $\zeta_r$
give the $16n$ block tail. A union bound over
the $n$ peers proves the source-cover statement.
\end{proof}

